\section{Descripción General}

\subsection{Perspectiva del Producto}
El sistema ``Gemelo Digital para Mantenimiento'' es un producto de software nuevo, diseñado como una solución complementaria e integradora dentro del ecosistema tecnológico de una planta industrial. No pretende reemplazar a los sistemas de gestión empresarial (ERP) o de gestión de activos (EAM/CMMS) existentes, sino actuar como una capa de visualización y operación avanzada que digitaliza los procesos que actualmente se realizan de forma manual o desconectada.

El sistema se estructura en una arquitectura distribuida que interactúa con las siguientes entidades externas:

\textbf{1. Interfaces del Sistema (Software):}
\begin{itemize}
    \item \textbf{Sistemas EAM/CMMS Existentes:} El producto se comunicará (vía API) con sistemas de registro central como Infor EAM o SAP PM para sincronizar el inventario de activos y el estado de las órdenes de trabajo, actuando como la interfaz de usuario moderna para estos sistemas ``backend''.
\end{itemize}

\textbf{2. Interfaces de Usuario:}
El sistema consta de dos componentes de interfaz principales, diseñados para contextos de uso específicos:
\begin{itemize}
    \item \textbf{Portal Web de Gestión (Escritorio):} Dirigido a Supervisores, Gerentes y Planeadores. Permite la visualización del Gemelo Digital en 3D, la gestión de tableros de control (dashboards), la configuración del sistema y el análisis de KPIs.
    \item \textbf{Aplicación Móvil (Tablet):} Dirigida a Técnicos de Mantenimiento e Inspectores HSEQ. Diseñada para operar en campo, permitiendo la visualización de rutas, consulta de manuales, escaneo de códigos QR para identificación de activos y reporte de actividades. \textbf{Esta interfaz debe contar con capacidad de operación desconectada (Offline-First)} para garantizar su funcionalidad en zonas de la planta sin cobertura de red.
\end{itemize}

\textbf{3. Interfaces de Hardware:}
\begin{itemize}
    \item \textbf{Dispositivos IoT y Sensores:} El sistema recibirá datos de telemetría (vibración, temperatura) provenientes de PLCs o microcontroladores (ej. ESP32, Raspberry Pi en entorno de prototipo) para alimentar el módulo de mantenimiento predictivo.
    \item \textbf{Dispositivos Móviles:} Interactuará con el hardware de las tablets (Cámara) para el escaneo de códigos QR de los activos y la captura de evidencia fotográfica de las reparaciones.
\end{itemize}

\subsection{Funciones del Producto}
El sistema se articula en cuatro módulos funcionales principales que cubren el ciclo de vida completo de la gestión de activos y mantenimiento:

\textbf{1. Gestión de Operaciones de Mantenimiento (MTTO)}: Funcionalidad central que permite la planificación, ejecución y seguimiento de las labores de mantenimiento. El sistema automatiza la generación de Órdenes de Trabajo (OT) preventivas y predictivas, gestiona el ciclo de vida de las solicitudes correctivas y facilita el cierre digital de órdenes en campo mediante dispositivos móviles. Incluye capacidades avanzadas para procesar informes de contratistas externos y estandarizar la entrada de datos técnicos.

\textbf{2. Gestión de Recursos y Activos (INV)}: Provee un repositorio centralizado y único para la ``Hoja de Vida'' de los activos, asegurando la trazabilidad de su historial, documentación técnica y jerarquía de componentes. Adicionalmente, gestiona el inventario de repuestos críticos (MRO), controlando las existencias y generando alertas automáticas de reabastecimiento para garantizar la disponibilidad de recursos durante las intervenciones.

\textbf{3. Visualización de Gemelo Digital (VIS)}: Ofrece una interfaz gráfica interactiva (2D/3D) que contextualiza la información operativa. Permite a los usuarios navegar espacialmente por la planta, visualizar capas de información técnica (rutas de aislamiento LOTO, permisos de trabajo) y acceder a vistas explosionadas de los equipos para identificar componentes sin necesidad de desarmarlos físicamente, mejorando la conciencia situacional y la seguridad.

\textbf{4. Inteligencia de Negocio y Administración (ADM)}: Módulo transversal encargado de transformar los datos operativos en información estratégica. Genera tableros de control (Dashboards) con indicadores clave de desempeño (MTBF, MTTR, Disponibilidad) para la toma de decisiones gerenciales. Asimismo, administra la seguridad del sistema, gestionando roles, usuarios y la interoperabilidad con plataformas externas.

\subsection{Características de los Usuarios}
La caracterización de los usuarios se ha derivado de la fase de elicitación de requisitos, específicamente de la entrevista con el experto del dominio y el análisis de procesos (BPA). Se han identificado tres categorías de usuarios con niveles de formación y habilidades técnicas diferenciadas. El diseño de la interfaz y la experiencia de usuario (UX) deben adaptarse a estas características para garantizar la adopción del sistema.

\textbf{1. Perfil Estratégico (Gerentes y Supervisores)}
\begin{itemize}
    \item \textbf{Roles:} Gerente de Planta, Supervisor de Mantenimiento.
    \item \textbf{Formación:} Profesional en Ingeniería (Mecánica, Industrial) o Administración. Frecuentemente con posgrados o certificaciones (CMRP).
    \item \textbf{Habilidades Técnicas:} Competentes en herramientas de ofimática y gestión (Excel avanzado, ERPs como SAP/Infor). No necesariamente expertos en hardware o configuración de sistemas.
    \item \textbf{Frecuencia de Uso:} Diaria (Supervisores) a Semanal (Gerentes).
    \item \textbf{Necesidad Principal:} Requieren interfaces tipo \textit{Dashboard} de lectura rápida, con datos sintetizados para la toma de decisiones. Valoran la eficiencia y la claridad sobre la profundidad técnica operativa.
\end{itemize}

\textbf{2. Perfil Operativo (Técnicos e Inspectores)}
\begin{itemize}
    \item \textbf{Roles:} Técnico de Mantenimiento, Inspector HSEQ, Contratista.
    \item \textbf{Formación:} Técnica o Tecnológica (una organización de talla mundial como el SENA o similar) en mecánica, electricidad o seguridad industrial.
    \item \textbf{Habilidades Técnicas:} Variable. Alta destreza en el manejo de maquinaria física, pero la familiaridad con software complejo puede ser limitada. Están acostumbrados al uso de smartphones, pero no necesariamente a PCs de escritorio.
    \item \textbf{Entorno de Uso:} Hostil (ruido, guantes de seguridad, luz solar directa, zonas sin conexión).
    \item \textbf{Necesidad Principal:} Requieren interfaces móviles (Tablet/Celular) con botones grandes, flujos lineales, mínima entrada de texto (uso de selectores y cámara) y capacidad de operar \textit{offline}.
\end{itemize}

\textbf{3. Perfil Técnico Especializado (Administración del Sistema)}
\begin{itemize}
    \item \textbf{Roles:} Ingeniero de Proyectos, Analista de Datos, Administrador del Sistema.
    \item \textbf{Formación:} Profesional en Ingeniería Mecatrónica, Sistemas o Electrónica.
    \item \textbf{Habilidades Técnicas:} Avanzadas. Conocimiento de CAD/CAM, bases de datos, y lógica de programación.
    \item \textbf{Frecuencia de Uso:} Diaria e intensiva.
    \item \textbf{Necesidad Principal:} Requieren acceso total a la configuración, capacidad para importar archivos complejos (modelos 3D), gestionar integraciones (APIs) y manipular grandes volúmenes de datos. La interfaz debe priorizar la funcionalidad y la densidad de información sobre la simplicidad.
\end{itemize}

\textbf{4. Asunciones sobre los Usuarios}
\begin{itemize}
    \item Se asume que todos los usuarios operativos cuentan con dispositivos móviles corporativos (Tablets) capaces de ejecutar la aplicación.
    \item Se asume que los roles estratégicos y especializados cuentan con estaciones de trabajo de escritorio con acceso a la red corporativa.
    \item Se asume que el \textbf{Analista de Datos} posee las competencias para interactuar directamente con los datos crudos o herramientas de BI externas, por lo que el sistema no requiere, en esta fase, una interfaz interna de desarrollo de modelos para este rol.
\end{itemize}

\subsection{Restricciones}
El diseño y la implementación del sistema están sujetos a las siguientes limitaciones obligatorias que deben ser respetadas durante el desarrollo:

\textbf{1. Restricciones de Hardware y Conectividad}
\begin{itemize}
    \item \textbf{Infraestructura de Red Limitada:} El sistema operará en un entorno industrial donde la conectividad a internet es intermitente en ciertas zonas. El software móvil \textbf{debe} ser capaz de operar en modo desconectado (\textit{Offline-first}) y sincronizar datos posteriormente.
    \item \textbf{Capacidad de Dispositivos:} La visualización del Gemelo Digital 3D debe ejecutarse en tablets de gama media estándar.
\end{itemize}

\textbf{2. Restricciones de Gestión de Proyecto y Temporales}
\begin{itemize}
    \item \textbf{Marco Temporal del Desarrollo:} El desarrollo del Producto Mínimo Viable (MVP) debe completarse dentro del marco temporal de la Etapa Lectiva, con una duración de \textbf{8 meses (Agosto 2026 a Abril 2027)}. Este plazo impone un límite estricto al alcance funcional.
\end{itemize}

\textbf{3. Estándares y Cumplimiento Normativo}
\begin{itemize}
    \item \textbf{Taxonomía de Activos (ISO 14224):} La estructura de la base de datos debe adherirse estrictamente al estándar \textbf{ISO 14224} para la recolección de datos de confiabilidad (ISO, 2016).
    \item \textbf{Formatos de Archivo:} El sistema solo permitirá la carga de informes externos en formato \textbf{PDF} y evidencias visuales en \textbf{JPG/PNG} por políticas de seguridad.
\end{itemize}

\textbf{4. Restricciones Económicas y de Licenciamiento}
\begin{itemize}
    \item \textbf{Costo de Licencias:} La solución debe basarse en tecnologías \textbf{Open Source} o de uso libre comercial para mitigar el costo de licencias de software propietario.
\end{itemize}

\textbf{5. Restricciones de Desarrollo}
\begin{itemize}
    \item \textbf{Curva de Aprendizaje y Tecnologías:} La solución debe desarrollarse utilizando lenguajes y frameworks con amplia documentación y soporte comunitario (ej. Java, JavaScript/TypeScript, Python) para asegurar la mantenibilidad y la viabilidad técnica por parte del equipo de desarrollo actual.
\end{itemize}

\subsection{Suposiciones y Dependencias}
Esta sección documenta los factores que se asumen como verdaderos y las dependencias externas que pueden afectar el alcance del proyecto:

\textbf{1. Suposiciones (Assumptions)}
\begin{itemize}
    \item \textbf{Recursos Humanos:} Se asume que el equipo de desarrollo mantendrá una composición mínima de 4 personas con disponibilidad de 6 horas diarias para las tareas asignadas.
    \item \textbf{Plataforma Base:} Se asume que los modelos 3D necesarios para el Gemelo Digital se obtendrán a través de la colaboración con el equipo de Ingeniería/CAD del cliente o se descargarán de catálogos de proveedores.
    \item \textbf{Ambiente de Producción:} Se asume que la empresa cliente cuenta con una infraestructura de red (Wi-Fi/LAN) operativa en sus oficinas administrativas.
\end{itemize}

\textbf{2. Dependencias (Dependencies)}
\begin{itemize}
    \item \textbf{Éxito de la Elicitación:} El alcance y la precisión de los Requisitos Específicos (Sección 3) dependen directamente de la calidad y el nivel de detalle de las respuestas obtenidas en la \textbf{Entrevista Semiestructurada con el experto de dominio}.
    \item \textbf{Fase de Evaluación (1.5 Meses):} La finalización del proyecto depende de la asignación exitosa del período de \textbf{1.5 meses de Evaluación Final} para realizar las Pruebas de Aceptación con el usuario y la documentación de cierre requerida para la certificación del SENA.
\end{itemize}

\subsection{Requisitos Futuros (Consideraciones para Versiones Posteriores)}
Las siguientes funcionalidades han sido identificadas como de alto valor, pero se excluyen del alcance del Producto Mínimo Viable (MVP) para concentrar el esfuerzo de desarrollo en el núcleo funcional del sistema. Se registrarán en el \textit{Product Backlog} para su consideración en futuras iteraciones.

\begin{itemize}
    \item \textbf{Integración Avanzada con Suites de Productividad:}
    \begin{itemize}
        \item \textbf{Descripción:} Sincronización nativa con plataformas como Google Workspace u Office 365 para la gestión de identidad (Single Sign-On), almacenamiento de documentos (Google Drive/SharePoint) y notificaciones avanzadas.
        \item \textbf{Justificación:} Aumentaría la adopción del usuario al integrarse con herramientas que ya utiliza diariamente.
    \end{itemize}
    \item \textbf{Módulo de Compras y Gestión Financiera:}
    \begin{itemize}
        \item \textbf{Descripción:} Funcionalidades para la emisión de órdenes de compra a proveedores, contabilidad de costos de mantenimiento y facturación.
        \item \textbf{Justificación:} Se excluye para limitar el alcance del módulo de inventario al control físico de cantidades, evitando la complejidad de la integración con sistemas financieros.
    \end{itemize}
    \item \textbf{Control Industrial Activo (SCADA/HMI):}
    \begin{itemize}
        \item \textbf{Descripción:} Capacidad para enviar comandos de control a la maquinaria (encender/apagar equipos, abrir/cerrar válvulas).
        \item \textbf{Justificación:} Se excluye explícitamente para garantizar la seguridad funcional de la planta, manteniendo el sistema en un modo de ``solo lectura'' respecto a la operación física.
    \end{itemize}
\end{itemize}
